\documentclass{homework}
\author{Tashfeen, Ahmad}
\class{CSCI 2114: Data Structures}
% \date{\today}
\title{Homework 2}
\address{Oklahoma City University, Petree College of Arts \& Sciences, Computer Science}

\usepackage{fontspec}
\setmonofont[Scale=MatchLowercase]{Iosevka SS07}

\begin{document} \maketitle

We start by reminding ourselves of the asymptotic order
notation\footnote{Sometimes also referred to as
  the \href{https://en.wikipedia.org/wiki/Big_O_notation}{Bachmann–Landau
    notation}.}, Let $f(x)$ and $g(x)$ be functions of a positive
$x$. We write,
\[
  f(x) = \O(g(x))
\]
when there is positive constant $c$ such that,
\[
  f(x) \leq cg(x)
\]
for all $x \geq n$. Function $f$ is then
stated as \textit{big oh} of $g$. Similarly, $f$ is \textit{big omega} of $g$ \ie
\[
  f(x) = \Omega(g(x))
\]
when there exist positive constants $c, n$ such that for all $x \geq n$ we have,
\[
  f(x) \geq cg(x)
\]
Lastly, if $f$ is both big-$\O$ and big-$\Omega$ of $g$,
\[
  \Omega(g(x)) = f(x) = \O(g(x))
\]
Then such a tight bound is stated as $f$ is big-$\Theta$ of $g$,
\[
  f(x) = \Theta(g(x))
\]
\question Consider the following Java subroutine,
\begin{lstlisting}[language=java]
public static void conditionalWord(int n) {
    for (int i = 0; i < n; i++)
        if (Math.random() < 0.5)
            taskA()
        else
            taskB()
}
\end{lstlisting}
\begin{enumerate}
  \item If on a certain machine the function \texttt{taskA()} takes
        2 seconds on each call in the loop and the function
        \texttt{taskB()} takes 4 seconds then for $n=10$, what is the
        total number of \textit{expected} seconds taken by the subroutine
        \texttt{public static void conditionalWord(int n)}?
  \item The subroutine \texttt{conditionalWord(int n)} is ran on a
        much faster machine reducing the time taken by \texttt{taskA()} on
        each call down to $\frac{1}{25}^\text{th}$ of a second and
        \texttt{taskB()} now takes $\frac{1}{5}^\text{th}$ of a
        second. For ${n=10}$, what is the new total number of
        seconds \textit{expected} by the subroutine \texttt{conditionalWord(int n)}?
  \item Assume that for any arbitrary $n$, \texttt{taskA()} takes
        $\log(n)$ steps while \texttt{taskB()} takes $2n^2$ steps. Let
        $T(n)$ be the total number of steps \textit{expected} by
        \texttt{conditionalWord(int n)}, what is $T(n)$?
  \item What is the \textit{best expected asymptotic} runtime
        complexity written as $T(n) = \Omega(g(n))$?
  \item What is the \textit{worst expected asymptotic} runtime
        complexity written as $T(n) = \O(g(n))$?
  \item If \texttt{taskB()} took $4\log(n)$ steps, what is the
        \textit{tight expected asymptotic} runtime complexity written
        as $T(n) = \Theta(g(n))$?
\end{enumerate}

\question In the code snippet bellow, we see two ways of getting the
tenth place digit of a Java \texttt{int} $n$. Give the runtime
complexity of each method using the \textit{big oh} notation. Justify
your answer.
\begin{lstlisting}[language=java]
System.out.println(n % 10); // the first way
System.out.println(String.valueOf(n).charAt(String.valueOf(n).length() - 1)); // the second way
\end{lstlisting}

\question Consider the following functions,
\begin{align*}
  a(x)      & = 2(x^3+1)(x^2-1)+2 \\
  b(x)      & = 2x^2              \\
  \alpha(x) & = x^2               \\
  \beta(x)  & = \pi x^2
\end{align*}
\begin{enumerate}
  \item Observe that $b(x) = \O(a(x))$ because $b(x) \leq a(x)$ past
        a certain $x=n$. What is the value of $n$ in this case?
  \item Observe that $\beta(x) = \O(\alpha(x))$ because there exists
        a $c$ such that $\beta(x) \leq c\alpha(x)$ for all $x > 0$. What
        is the value of $c$ in this case?
\end{enumerate}

\question Given bellow is a Java subroutine that sorts an array of
non-negative Java integers in ascending order.
\begin{lstlisting}[language=java]
public static void sortIntegers(int[] toSort) {
    int i = 0, j = 0, k = 0, max = Integer.MIN_VALUE;
    for (i = 0; i < toSort.length; i++)
        max = toSort[i] > max ? toSort[i] : max;
    int[] counts = new int[max + 1];
    for (i = 0; i < toSort.length; i++)
        counts[toSort[i]]++;
    for (i = 0; i < counts.length; i++)
        for (j = 0; j < counts[i]; j++)
            toSort[k++] = i;
}
\end{lstlisting}

Assume that the length of the input parameter \texttt{int[]} is $n$
\ie \texttt{toSort.length} $= n$. Let $T(n)$ be the worst and average
case complexity of the algorithm's runtime and $S(\texttt{toSort})$ be the worst
case complexity of memory-space.

\begin{enumerate}
  \item What are $T(n), S(\texttt{toSort})$?
  \item What are $\O(T(n)), \O(S(\texttt{toSort}))$?
  \item Look up and state the \textit{big}-$\O$ of the average case
        complexity of Java's built-in
        \href{https://tinyurl.com/jw75mfa}{\texttt{Arrays.sort(int[])}}. Is
        this better than $\O(T(n))$ from the previous step of this question?
\end{enumerate}

\question Given bellow is the Java implementation of the sieve of
Eratosthenes. This particular implementation marks all the composite
numbers as \texttt{true} since Java allocates \texttt{false} to all
the elements of a newly created boolean arrays.
\begin{lstlisting}[language=java]
public static void eratosthenes(boolean[] toSieve) {
    toSieve[0] = true;
    toSieve[1] = true;
    for (int i = 2; i < Math.sqrt(toSieve.length); i++)
        if (!toSieve[i])
            for (int j = i*i; j < toSieve.length; j += i)
                toSieve[j] = true;
}
\end{lstlisting}
Give an upper-bound on the runtime complexity of this particular
implementation of the sieve of Eratosthenes. The better upper-bound
you give, the more credit you get. Justify your answer.
\section*{Submission Instructions}
Submit a PDF file with your answers.
\end{document}
